% !TeX program = xelatex
%
% IMPE LaTeX System — Canonical Showcase (v1.1.0)
% Build with scripts/build_showcase.ps1 from the repository root.

\documentclass[11pt,twoside]{impeart}

\HeaderTitle{IMPE Canonical Showcase}

\UseTemplateSet{
  fonts = {libertinus,shanggu},
  features = {
    math,
    tables,
    image,
    lists_envs,
    citations,
    index,
    hyperlinks,
    headers
  }
}

\setmonofont{Latin Modern Mono}

% Local families used by the current specimens and style matrix.
\UseFonts{
  noto, times, gentium, charis, mlmodern, anatolian,
  bopomofo, cuneiform, italic, hungarian, runic,
  hindi, sanskrit, devanagari, brahmi, tibetan,
  arabic, urdu, aramaic, nabataean, hebrew,
  syriac, syriac_eastern, kharosthi, egyptian, pahlavi_parthian, pahlavi_inscriptional,
  pahlavi_psalter, avestan, manichaean, phoenician, samaritan,
  sogdian, sogdian_old, chinese_simplified, chinese_traditional, japanese,
  wenjin, korean, tangut, khitan_small, mongolian,
  mongolian_baiti, manchu, segoe, thai, turkic,
  uyghur, vietnamese_quocngu, vietnamese_hannom
}[local]

% Automatic Unicode-range routing for the current mixed-script specimen.
\UseFonts{armenian,georgian,coptic,glagolitic,tamil}

\addbibresource{references.bib}

\hypersetup{
  pdftitle   = {IMPE LaTeX System — Canonical Showcase},
  pdfauthor  = {Kelvin Yang},
  pdfsubject = {Multilingual and scholarly typesetting}
}

\setcounter{tocdepth}{2}
\raggedbottom

\newcommand{\ShowcaseStyles}[2]{%
  \leavevmode #1{#2} &
  \leavevmode #1{\textbf{#2}} &
  \leavevmode #1{\textit{#2}} &
  \leavevmode #1{\textbf{\textit{#2}}} &
  \leavevmode #1{\textsf{#2}} &
  \leavevmode #1{\textsf{\textbf{#2}}} &
  \leavevmode #1{\textsf{\textit{#2}}} &
  \leavevmode #1{\textsf{\textbf{\textit{#2}}}}%
}

\newcommand{\StyleRow}[3]{%
  #1 &
  \ShowcaseStyles{#2}{#3} \\
}

\newcommand{\StyleCategory}[1]{%
  \midrule
  \multicolumn{9}{l}{\textbf{#1}}\\
  \midrule
}

\title{IMPE \LaTeX~System}
\subtitle{Canonical Typesetting Showcase}
\author{Kelvin Yang}
\date{v1.1.0}

\begin{document}

\maketitle

\begin{abstract}
This document is the canonical visual specimen for the current IMPE \LaTeX~System public interface. It combines a local-family style audit, a deliberately
small Unicode-range routing regression, dense multilingual paragraph
composition with explicit local overrides, right-to-left text, vertical
writing, generic complex shaping, CJK and Han family selection, and the
principal document features provided by IMPE. It is intended both as a public
showcase and as a release-time visual regression specimen.\footnote{The
showcase is intentionally broader than the focused regression fixtures under
\texttt{tests/}; those fixtures isolate individual failures.}
\end{abstract}

\tableofcontents



% ==========================================================================
\section{Local Font Families and Styles}
% ==========================================================================

The following typeface atlas pairs concise script specimens with eight
corresponding text styles. Repeated specimens make differences between
typefaces directly comparable; where a face is not supplied, the registered
family's normal fallback applies. Egyptian format controls and Khitan clusters
use their native two-dimensional layouts; Inner Asian samples use vertical
writing. The matrix continues across landscape pages as needed.

\afterpage{%
\begin{landscape}
\small
\setlength{\tabcolsep}{4pt}

\begin{longtable}{
  @{}
  P{2.80cm}
  *{8}{M{2.20cm}}
  @{}
}
\caption{Typeface specimens in eight styles.}\label{tab:typeface-atlas}\\
\toprule
Typeface &
Regular &
Bold &
Italic &
Bold italic &
Sans &
Sans bold &
Sans italic &
Sans bold italic \\
\midrule
\endfirsthead

\toprule
Typeface &
Regular &
Bold &
Italic &
Bold italic &
Sans &
Sans bold &
Sans italic &
Sans bold italic \\
\midrule
\endhead

\StyleCategory{Latin typefaces}
\StyleRow{CMU Serif}{\Writing{[cmu]latin}}{Latina}
\StyleRow{Noto}{\NOT}{Latina}
\StyleRow{Times}{\TIM}{Latina}
\StyleRow{Gentium Plus}{\GEN}{Latina}
\StyleRow{Charis SIL}{\CHA}{Latina}
\StyleRow{Libertinus}{\LIB}{Latina}
\StyleRow{MLModern}{\MLM}{Latina}

\StyleCategory{European and Caucasian scripts}
\StyleRow{Carian}{\CA}{𐊨𐊣𐊠𐊦𐊹𐊸}
\StyleRow{Coptic}{\CO}{ⲛⲟⲩⲧⲉ}
\StyleRow{Glagolitic}{\GL}{ⰳⰾⰰⰳⱁⰾ}
\StyleRow{Old Italic}{\OI}{𐌌𐌉 𐌔𐌖𐌈𐌉}
\StyleRow{Old Hungarian}{\OH}{𐳢𐳛𐳮𐳁𐳤}
\StyleRow{Runic}{\RU}{ᚠᚢᚦᚨᚱᚲ}
\StyleRow{Armenian}{\HY}{հայերեն}
\StyleRow{Georgian}{\KA}{ქართული}

\StyleCategory{South Asian and Tibetan scripts}
\StyleRow{Devanagari · Hindi}{\HI}{देवनागरी}
\StyleRow{Devanagari · Sanskrit}{\SA}{देवनागरी}
\StyleRow{Devanagari}{\DEV}{देवनागरी}
\StyleRow{Tamil}{\TA}{தமிழ்}
\StyleRow{Brahmi}{\BR}{𑀥𑀫𑁆𑀫}
\StyleRow{Brahmi · Segoe}{\SEG}{𑀥𑀫𑁆𑀫}
\StyleRow{Kharoshthi}{\KH}{𐨢𐨿𐨪𐨨}
\StyleRow{Tibetan}{\TI}{བོད་ཡིག}
\StyleRow{Thai}{\TH}{ภาษาไทย}

\StyleCategory{West Asian and Iranian scripts}
\StyleRow{Arabic}{\AR}{العربية}
\StyleRow{Arabic · Nastaliq}{\UR}{نستعلیق}
\StyleRow{Imperial Aramaic}{\IA}{𐡀𐡓𐡌𐡉𐡀}
\StyleRow{Nabataean}{\NB}{𐢕𐢃𐢋𐢈}
\StyleRow{Hebrew}{\HE}{עברית}
\StyleRow{Syriac}{\SY}{ܣܘܪܝܝܐ}
\StyleRow{Syriac · Eastern}{\SYE}{ܣܘܪܝܝܐ}
\StyleRow{Inscriptional Parthian}{\PAR}{𐭃𐭐𐭉𐭅𐭓}
\StyleRow{Inscriptional Pahlavi}{\PAH}{𐭣𐭯𐭩𐭥𐭥}
\StyleRow{Psalter Pahlavi}{\PSP}{𐮃𐮎𐮈𐮅𐮅}
\StyleRow{Avestan}{\AV}{𐬀𐬴𐬆𐬨}
\StyleRow{Manichaean}{\MA}{𐫖𐫀𐫗𐫏}
\StyleRow{Phoenician}{\PH}{𐤁𐤏𐤋}
\StyleRow{Samaritan}{\SM}{ࠔࠋࠅࠌ}
\StyleRow{Sogdian}{\SG}{𐼼𐼴𐼶𐼹𐼷𐼸}
\StyleRow{Old Sogdian}{\SGO}{𐼑𐼇𐼄𐼌𐼊𐼋}

\StyleCategory{Ancient Near Eastern scripts}
\StyleRow{Cuneiform}{\CU}{𒁾𒊬}
\StyleRow{Egyptian Hieroglyphs}{\EG}{𓆓𓐳𓂧𓇋𓐰𓈖𓌞𓅱𓀀𓇋𓈎𓐰𓂋𓐰𓏛}

\StyleCategory{East Asian scripts and Han families}
\StyleRow{Bopomofo}{\ZY}{ㄅㄆㄇㄈ}
\StyleRow{Han · Simplified}{\SC}{汉字}
\StyleRow{Han · Traditional}{\TC}{漢字}
\StyleRow{Han · Hiragana · Katakana}{\JP}{漢字かなカナ}
\StyleRow{Han · WenJin}{\WJ}{漢字}
\StyleRow{Han · Shanggu}{\Writing{[shanggu]chinese}}{漢字}
\StyleRow{Han · Sim}{\Writing{[simsun]chinese}}{汉字}
\StyleRow{Hangul · Han}{\KR}{한글漢字}
\StyleRow{Tangut}{\TG}{𗼇𘝞}

\StyleCategory{Inner Asian scripts}
\StyleRow{Khitan Small Script}{\KHS}{𘰺𘱹𘰻 𘮦𖿤𘲲}
\StyleRow{Mongolian}{\MOv}{ᠮᠣᠩᠭᠣᠯ}
\StyleRow{Mongolian Baiti}{\MObv}{ᠮᠣᠩᠭᠣᠯ}
\StyleRow{Manchu}{\MCv}{ᠮᠠᠨᠵᡠ}
\StyleRow{Old Turkic}{\OT}{𐱅𐰇𐰼𐰜}
\StyleRow{Old Uyghur}{\UYv}{𐽳𐽶𐽲𐽳𐽾}

\StyleCategory{Vietnamese writing traditions}
\StyleRow{Latin · Quốc Ngữ}{\VI}{Quốc ngữ}
\StyleRow{Han · Chữ Nôm}{\HN}{𡨸喃}

\bottomrule
\end{longtable}

\end{landscape}
}

% ==========================================================================
\section{Multilingual and Mixed-Script Composition}
% ==========================================================================

Written traditions carry voices across languages, centuries, and regions.
The passages here bring together reflections on books and learning alongside
short proverbial and scriptural lines, composed in their own writing systems.

\subsection{Polyglot Composition}

The same page can accommodate many scripts without assigning any single one
the role of a translation or a transcription. In the following continuous
paragraph, five writing systems are set directly in Unicode, while the others
retain their chosen local typefaces.\index{Multilingual composition}
\index{Unicode-range routing}

\noindent
Գրքերը մեզ պատմում են անցյալի մարդկանց մասին։
ძველ წიგნებში მრავალი თაობის აზრი და გამოცდილებაა შემონახული.
Ⲟⲩⲟϩ ⲡⲓⲗⲟⲅⲟⲥ ⲁϥϣⲱⲡⲓ ⲛ̀ⲟⲩⲥⲁⲣⲝ ⲟⲩⲟϩ ⲁϥϣⲱⲡⲓ ⲛ̀ϧⲏⲧⲉⲛ.
Ⰲⱏ ⱀⰰⱍⰰⰾⱑ ⰱⱑⰰⱎⰵ ⱄⰾⱁⰲⱁ, ⰺ ⱄⰾⱁⰲⱁ ⰱⱑⰰⱎⰵ ⰽⱏ ⰱⱁⰳⱆ.
\SA{विद्या ददाति विनयं।}
கற்க கசடறக் கற்பவை கற்றபின் நிற்க அதற்குத் தக.
\TI{ཤེས་རབ་དང་སྙིང་རྗེ་ནི་མིའི་ཆོས་ཀྱི་རྩ་བ་ཡིན།}
\AR{الكتاب خير جليس، والعلم نور.}
\UR{کتابیں ہمیں ماضی کی دنیا سے ملاتی ہیں۔}
\HE{החכמה טובה מכסף ומזהב.}
\SY{ܒܪܫܝܬ ܐܝܬܘܗܝ ܗܘܐ ܡܠܬܐ.}
\SYE{ܒܪܫܝܬ ܐܝܬܘܗܝ ܗܘܐ ܡܠܬܐ.}
\SEG{𑀲𑀢𑁆𑀬𑀁 𑀯𑀤, 𑀥𑀭𑁆𑀫𑀁 𑀘𑀭.}
\JP{古い図書館でカタログをめくり、昔の人の手紙を読みました。}
\KR{古典을 읽으면 오늘의 질문을 다시 생각하게 됩니다.}
\TC{舊書中的一句話，有時會在數百年後重新被人理解。}

\subsection{Mixed Writing Directions}

A page may also alternate between horizontal sentences and vertically
written words. Here the Mongolian and Manchu passages rise between ordinary
horizontal lines, while each writing system retains its own orientation.

\noindent
\TC{翻開一本舊書，紙頁上的文字把我們帶到另一個時代。}
\MOv{ᠮᠣᠩᠭᠣᠯ ᠬᠡᠯᠡ ᠪᠣᠯ ᠮᠣᠩᠭᠣᠯᠴᠤᠳ ‍ᠤᠨ ᠬᠡᠯᠡ ᠮᠥᠨ᠃}
\JP{古い手紙を読むと、文字の形から書き手の暮らしまで想像できます。}
\MCv{ᠪᡳ ᠮᠠᠨᠵᡠ ᡤᡳᠰᡠᠨ ᠪᡝ ᡨᠠᠴᡳᠮᠪᡳ᠂ ᡝᡵᡝ ᠪᡳᡨᡥᡝ ᠪᡝ ᡥᡡᠯᠠᠮᠪᡳ᠃}
\KR{오래된 기록을 읽는 일은 낯선 사람의 목소리에 귀를 기울이는 일입니다.}
\TC{文字的方向雖然不同，閱讀仍讓相隔久遠的人彼此相遇。}


% ==========================================================================
\section{Right-to-Left Scripts}
% ==========================================================================

Arabic, Urdu, Hebrew, and Syriac begin their lines at the right-hand edge,
yet their books also accommodate dates, shelfmarks, and occasional passages in
other writing systems. These pages bring together modern prose, verse, and
older scriptural traditions.\index{Right-to-left text}

\subsection{Arabic}

In a quiet manuscript room, a catalogue entry may be the beginning of a
much longer story.

\AR{%
في قاعة المخطوطات رفوف من الخشب الداكن، وعلى الرفوف كتب نُسخت في مدن مختلفة وعصور متباعدة. تحمل بعض الصفحات أسماء ناسخين وقرّاء، وتحفظ الهوامش أسئلة صغيرة تركها أشخاص لا نعرف عن حياتهم إلا القليل. تُظهر إحدى النسخ آثار التصحيحات التي أضافها ناسخ لاحق، وتروي حواشيها كيف انتقل الكتاب بين أيدي القرّاء. وحين يقارن الباحث نسختين من الكتاب نفسه يجد أن اختلاف كلمة واحدة قد يكشف طريق النص من مدينة إلى أخرى. لهذا لا تكون القراءة عودة إلى الماضي وحده؛ إنها حوار جديد مع أناس حفظوا الكلمات ونقلوها إلى من جاء بعدهم.
\par
}

A fully vocalized couplet by al-Mutanabbī offers another typographic
rhythm, with its vowel marks and joined letterforms:

\AR{أَعَزُّ مَكَانٍ فِي الدُّنَى سَرْجُ سَابِحٍ، وَخَيْرُ جَلِيسٍ فِي الزَّمَانِ كِتَابُ.}

\subsection{Urdu}

\UR{%
ایک پرانے کتب خانے میں داخل ہوتے ہی کاغذ اور روشنائی کی ایک مانوس خوشبو محسوس ہوتی ہے۔ لکڑی کی الماریوں میں رکھی ہوئی کتابوں پر زمانے کے نشان صاف دکھائی دیتے ہیں، لیکن ان کے صفحات آج بھی نئے سوال پیدا کرتے ہیں۔ بعض نسخوں کے حاشیوں میں پچھلے قارئین نے اپنے خیالات لکھے ہیں؛ کہیں کسی مشکل لفظ کا مطلب درج ہے اور کہیں صرف ایک چھوٹا سا سوالیہ نشان۔ سن ۱۹۲۷ کی ایک یادداشت میں کتاب کے مالک نے لکھا تھا کہ اسے یہ نسخہ سفر کے دوران ملا۔ اب وہی کتاب دوسرے شہر کی لائبریری میں موجود ہے، جہاں ہر نیا قاری اس کی کہانی میں اپنا حصہ شامل کرتا ہے۔ یوں ایک کتاب صرف الفاظ کا مجموعہ نہیں رہتی بلکہ مختلف زمانوں اور لوگوں کے درمیان رابطے کا ذریعہ بن جاتی ہے۔
\par
}

\subsection{Hebrew}

\HE{%
כשנכנסים לספרייה ישנה, אפשר לפעמים לשמוע את רשרוש הדפים יותר מכל קול אחר. על שולחן העץ מונחים כתבי יד, מהדורות מודפסות ומחברת של קורא שכבר איננו. בשוליו של אחד הספרים נכתבה הערה קצרה בשנת 1896, ולידה סימן שאלה קטן. כעבור דורות פותח קורא אחר את אותו העמוד וממשיך את השיחה שהתחילה לפניו. כך נשמרות בספרייה לא רק המילים שבספרים, אלא גם עקבותיהם של האנשים שקראו אותם, העירו עליהם והעבירו אותם לידיים אחרות.
\par
}

The opening of Genesis presents a different Hebrew page: fully pointed
consonants, cantillation marks, and the verse-ending punctuation.

\HE{%
בְּרֵאשִׁ֖ית בָּרָ֣א אֱלֹהִ֑ים אֵ֥ת הַשָּׁמַ֖יִם וְאֵ֥ת הָאָֽרֶץ׃
\par
}

\subsection{Syriac}

The opening of the Gospel of John in the Syriac Peshitta places words
and light at the heart of its narrative. The following continuous passage
is taken from John 1:2--5.

\SY{%
ܗܢܐ ܐܝܬܘܗܝ ܗܘܐ ܒܪܫܝܬ ܠܘܬ ܐܠܗܐ ܀ ܟܠ ܒܐܝܕܗ ܗܘܐ ܘܒܠܥܕܘܗܝ ܐܦܠܐ ܚܕܐ ܗܘܬ ܡܕܡ ܕܗܘܐ ܀ ܒܗ ܚܝܐ ܗܘܐ ܘܚܝܐ ܐܝܬܝܗܘܢ ܢܘܗܪܐ ܕܒܢܝܢܫܐ ܀ ܘܗܘ ܢܘܗܪܐ ܒܚܫܘܟܐ ܡܢܗܪ ܘܚܫܘܟܐ ܠܐ ܐܕܪܟܗ ܀
\par
}

Two traditional typefaces set the same opening verse, allowing their
letterforms to be read alongside one another.

\noindent
\begin{minipage}[t]{0.47\linewidth}
\textbf{General Syriac}

\medskip

\SY{%
ܒܪܫܝܬ ܐܝܬܘܗܝ ܗܘܐ ܡܠܬܐ ܘܗܘ ܡܠܬܐ ܐܝܬܘܗܝ ܗܘܐ ܠܘܬ ܐܠܗܐ ܘܐܠܗܐ ܐܝܬܘܗܝ ܗܘܐ ܗܘ ܡܠܬܐ ܀
\par
}
\end{minipage}
\hfill
\begin{minipage}[t]{0.47\linewidth}
\textbf{Eastern Syriac}

\medskip

\SYE{%
ܒܪܫܝܬ ܐܝܬܘܗܝ ܗܘܐ ܡܠܬܐ ܘܗܘ ܡܠܬܐ ܐܝܬܘܗܝ ܗܘܐ ܠܘܬ ܐܠܗܐ ܘܐܠܗܐ ܐܝܬܘܗܝ ܗܘܐ ܗܘ ܡܠܬܐ ܀
\par
}
\end{minipage}

% ==========================================================================
\section{Vertical Writing}
% ==========================================================================

Mongolian, Manchu, and Old Uyghur are written in vertical lines. In a
horizontal IMPE document, local vertical runs retain that orientation while
individual word-sized units can wrap within the surrounding page. The
following passages compare the two settings and their different letterforms.
\index{Vertical writing}

\subsection{Mongolian}

The same Mongolian passage appears in horizontal and vertical settings.

\noindent
\begin{minipage}[t]{0.47\linewidth}
\raggedright
\textbf{Horizontal reference}

\medskip

\MO{ᠮᠣᠩᠭᠣᠯ ᠬᠡᠯᠡ ᠪᠣᠯ ᠮᠣᠩᠭᠣᠯᠴᠤᠳ ‍ᠤᠨ ᠬᠡᠯᠡ ᠮᠥᠨ᠃ ᠮᠣᠩᠭᠣᠯ ᠪᠢᠴᠢᠭ ᠨᠢ ᠪᠣᠰᠤᠭ᠎ᠠ ᠴᠢᠭᠯᠡᠯ ᠢᠶᠡᠷ ᠪᠢᠴᠢᠭᠳᠡᠨ᠎ᠡ᠃}
\end{minipage}
\hfill
\begin{minipage}[t]{0.47\linewidth}
\raggedright
\textbf{Vertical}

\medskip

\MOv{ᠮᠣᠩᠭᠣᠯ ᠬᠡᠯᠡ ᠪᠣᠯ ᠮᠣᠩᠭᠣᠯᠴᠤᠳ ‍ᠤᠨ ᠬᠡᠯᠡ ᠮᠥᠨ᠃ ᠮᠣᠩᠭᠣᠯ ᠪᠢᠴᠢᠭ ᠨᠢ ᠪᠣᠰᠤᠭ᠎ᠠ ᠴᠢᠭᠯᠡᠯ ᠢᠶᠡᠷ ᠪᠢᠴᠢᠭᠳᠡᠨ᠎ᠡ᠃}
\end{minipage}

\medskip

The regular Mongolian family and Mongolian Baiti give the same short
phrase different appearances in vertical setting:

\noindent
\begin{minipage}[t]{0.47\linewidth}
\raggedright
\textbf{Mongolian}

\medskip

\MOv{ᠮᠣᠩᠭᠣᠯ ᠪᠢᠴᠢᠭ ᠨᠢ ᠪᠣᠰᠤᠭ᠎ᠠ}
\end{minipage}
\hfill
\begin{minipage}[t]{0.47\linewidth}
\raggedright
\textbf{Mongolian Baiti}

\medskip

\MObv{ᠮᠣᠩᠭᠣᠯ ᠪᠢᠴᠢᠭ ᠨᠢ ᠪᠣᠰᠤᠭ᠎ᠠ}
\end{minipage}

\bigskip

\subsection{Manchu}

The same Manchu passage appears in both orientations, allowing direct
comparison without changing the text.

\noindent
\begin{minipage}[t]{0.47\linewidth}
\raggedright
\textbf{Horizontal reference}

\medskip

\MC{ᠮᠠᠨᠵᡠ ᡤᡳᠰᡠᠨ ᠪᡝ ᠮᠠᠨᠵᡠ ᠪᡳᡨᡥᡝ ‍ᡳ ᠠᡵᠠᠮᠪᡳ᠃ ᠪᡳ ᡝᡵᡝ ᠪᡳᡨᡥᡝ ᡩᡝ ᠯᠣᠨᡤᡤᠣ ᡤᡳᠰᡠᠨ ᠪᡝ ᠠᡵᠠᠮᠪᡳ᠂ ᡨᡝᡵᡝᠴᡳ ᡤᡳᠰᡠᠨ ᡨᠠᡨᠠᠮᡝ ᡤᡝᠨᡝᠮᠪᡳ᠃}
\end{minipage}
\hfill
\begin{minipage}[t]{0.47\linewidth}
\raggedright
\textbf{Vertical}

\medskip

\MCv{ᠮᠠᠨᠵᡠ ᡤᡳᠰᡠᠨ ᠪᡝ ᠮᠠᠨᠵᡠ ᠪᡳᡨᡥᡝ ‍ᡳ ᠠᡵᠠᠮᠪᡳ᠃ ᠪᡳ ᡝᡵᡝ ᠪᡳᡨᡥᡝ ᡩᡝ ᠯᠣᠨᡤᡤᠣ ᡤᡳᠰᡠᠨ ᠪᡝ ᠠᡵᠠᠮᠪᡳ᠂ ᡨᡝᡵᡝᠴᡳ ᡤᡳᠰᡠᠨ ᡨᠠᡨᠠᠮᡝ ᡤᡝᠨᡝᠮᠪᡳ᠃}
\end{minipage}

\bigskip

\subsection{Old Uyghur}

A provisional Unicode rendering based on the published transliteration
of the colophon of the Turfan manuscript U~338 supplies a continuous
Old Uyghur passage. It is presented as a typographic specimen rather
than a critical edition.

\noindent
\begin{minipage}[t]{0.47\linewidth}
\raggedright
\textbf{Horizontal reference}

\medskip

\UY{𐽰𐽳𐽸 𐽶𐽶𐾁 𐽾𐽰𐽹 𐽰𐽰𐽶 𐽰𐽳𐽶𐽽 𐽰𐽳𐾀𐽳𐽴𐽲𐽰 𐽹𐽺 𐽼𐽰𐽷𐽳𐽴 𐽼𐽶𐾀𐽶𐾀𐽶𐽹 𐽼𐽳𐽶𐽰𐽺𐽶 𐽰𐾀𐽹 𐽶𐽳𐽺𐽰𐽺𐽲𐽰 𐾀𐽰𐽷𐽻𐽳𐽺 𐽰𐽰𐽹𐽶𐽺 𐽰𐽰𐽹𐽶𐽺 𐾀𐽶𐽼 𐽹𐽳𐽺𐽶 𐾀𐽰𐽷 𐽼𐽶𐽶𐽻 𐽲𐽰𐾀𐽰 𐽼𐽶𐾀𐽶𐽼 𐽰𐽱𐽸𐽰 𐽰𐽳𐽲𐽳𐾁 𐽲𐽶𐽴 𐽲𐽰 𐽰𐽳𐽲𐽶𐾀𐽶𐽼 𐽰𐽻𐽶𐽸𐽹𐽰𐽷𐽶 𐽼𐽳𐾁𐽴𐽳𐽺 𐽰𐽹𐽶𐽺}
\end{minipage}
\hfill
\begin{minipage}[t]{0.47\linewidth}
\raggedright
\textbf{Vertical}

\medskip

\UYv{𐽰𐽳𐽸 𐽶𐽶𐾁 𐽾𐽰𐽹 𐽰𐽰𐽶 𐽰𐽳𐽶𐽽 𐽰𐽳𐾀𐽳𐽴𐽲𐽰 𐽹𐽺 𐽼𐽰𐽷𐽳𐽴 𐽼𐽶𐾀𐽶𐾀𐽶𐽹 𐽼𐽳𐽶𐽰𐽺𐽶 𐽰𐾀𐽹 𐽶𐽳𐽺𐽰𐽺𐽲𐽰 𐾀𐽰𐽷𐽻𐽳𐽺 𐽰𐽰𐽹𐽶𐽺 𐽰𐽰𐽹𐽶𐽺 𐾀𐽶𐽼 𐽹𐽳𐽺𐽶 𐾀𐽰𐽷 𐽼𐽶𐽶𐽻 𐽲𐽰𐾀𐽰 𐽼𐽶𐾀𐽶𐽼 𐽰𐽱𐽸𐽰 𐽰𐽳𐽲𐽳𐾁 𐽲𐽶𐽴 𐽲𐽰 𐽰𐽳𐽲𐽶𐾀𐽶𐽼 𐽰𐽻𐽶𐽸𐽹𐽰𐽷𐽶 𐽼𐽳𐾁𐽴𐽳𐽺 𐽰𐽹𐽶𐽺}
\end{minipage}


% ==========================================================================
\section{Complex Shaping and Stacked Sequences}
% ==========================================================================

Two kinds of composition meet on these pages. Devanagari, Tibetan,
Brahmi, Kharoshthi, and Thai rely on the shaping of ordinary Unicode text
by the selected OpenType fonts. Egyptian Hieroglyphs and Khitan Small
Script instead use XSR to arrange signs into two-dimensional units while
preserving their original Unicode sequences.
\index{OpenType shaping}\index{Two-dimensional composition}

\subsection{Devanagari}

Devanagari makes conjunct consonants, vowel signs, and the virama visible
in the shapes of its syllables. The Sanskrit family sets the following
forms with the same typeface used in the verse below.

\begin{TableBookX}
  {L L}
  {Devanagari shaping specimens}
  {tab:deva-shaping}
Sequence & Specimen \\
\midrule
kṣa & \SA{क्ष} \\
jña & \SA{ज्ञ} \\
śra & \SA{श्र} \\
stra & \SA{स्त्र} \\
ddhya & \SA{द्ध्य} \\
dharma & \SA{धर्म} \\
dharmakṣetre & \SA{धर्मक्षेत्रे} \\
saṃskṛtam & \SA{संस्कृतम्} \\
oṃ & \SA{ॐ} \\
\end{TableBookX}

The opening of the \emph{Bhagavadgītā} (1.1):

\begin{quote}
\upshape
\SA{धृतराष्ट्र उवाच।
धर्मक्षेत्रे कुरुक्षेत्रे समवेता युयुत्सवः।
मामकाः पाण्डवाश्चैव किमकुर्वत सञ्जय॥}
\end{quote}

\subsection{Tibetan}

Tibetan syllables can contain stacked consonants within a single
\emph{tsheg}-delimited unit. The short forms below are followed by a
connected passage from the traditional Four Immeasurables prayer.

\begin{TableBookX}
  {L L}
  {Tibetan consonant-stack specimens}
  {tab:tibetan-shaping}
Transliteration & Specimen \\
\midrule
bsgrubs & \TI{བསྒྲུབས་} \\
brgyud & \TI{བརྒྱུད་} \\
bstan & \TI{བསྟན་} \\
rgyal & \TI{རྒྱལ་} \\
dbyangs & \TI{དབྱངས་} \\
\end{TableBookX}

\TI{སེམས་ཅན་ཐམས་ཅད་བདེ་བ་དང་བདེ་བའི་རྒྱུ་དང་ལྡན་པར་གྱུར་ཅིག །
སེམས་ཅན་ཐམས་ཅད་སྡུག་བསྔལ་དང་སྡུག་བསྔལ་གྱི་རྒྱུ་དང་བྲལ་བར་གྱུར་ཅིག །
སེམས་ཅན་ཐམས་ཅད་སྡུག་བསྔལ་མེད་པའི་བདེ་བ་དང་མི་འབྲལ་བར་གྱུར་ཅིག །
སེམས་ཅན་ཐམས་ཅད་ཉེ་རིང་ཆགས་སྡང་དང་བྲལ་བའི་བཏང་སྙོམ་ལ་གནས་པར་གྱུར་ཅིག །}

\medskip
\noindent\textbf{Sanskrit transcribed in Tibetan letters.}
The \emph{Heart Sūtra} mantra illustrates Sanskrit consonant stacks,
long-vowel marks, anusvāra, and the \emph{svāhā} cluster without
switching away from the Tibetan-script font:

\begin{quote}
\upshape
\TI{ཏདྱ་ཐཱ། ག་ཏེ་ག་ཏེ་པཱ་ར་ག་ཏེ།
པཱ་ར་སཾ་ག་ཏེ། བོ་དྷི་སྭཱ་ཧཱ།}
\end{quote}

\subsection{Brahmi, Kharoshthi, and Thai}

The same Sanskrit sentences from the \emph{Taittirīya Upaniṣad}
(1.11.1: \emph{satyaṃ vada, dharmaṃ cara, svādhyāyān mā pramadaḥ})
appear in two Brahmi fonts, giving their conjuncts and diacritics a
direct side-by-side comparison. The Kharoshthi specimen represents
those same Sanskrit sentences in Kharoshthi characters, in
right-to-left order. These are illustrative script transcriptions,
not claims about the script of the surviving Upaniṣad manuscript.
The Thai passage is continuous contemporary prose, allowing
word-level line breaking and vowel/tone placement to be inspected
within a constrained table column.

\begin{TableBookX}
  {P{0.26\linewidth} L}
  {Indic and Thai composition in continuous text}
  {tab:generic-shaping}
Script / family & Specimen \\
\midrule
Brahmi / Noto &
  \BR{𑀲𑀢𑁆𑀬𑀁 𑀯𑀤𑁇 𑀥𑀭𑁆𑀫𑀁 𑀘𑀭𑁇 𑀲𑁆𑀯𑀸𑀥𑁆𑀬𑀸𑀬𑀸𑀦𑁆 𑀫𑀸 𑀧𑁆𑀭𑀫𑀤𑀂𑁇} \\
Brahmi / Segoe Historic &
  \SEG{𑀲𑀢𑁆𑀬𑀁 𑀯𑀤𑁇 𑀥𑀭𑁆𑀫𑀁 𑀘𑀭𑁇 𑀲𑁆𑀯𑀸𑀥𑁆𑀬𑀸𑀬𑀸𑀦𑁆 𑀫𑀸 𑀧𑁆𑀭𑀫𑀤𑀂𑁇} \\
Kharoshthi &
  \KH{𐨯𐨟𐨿𐨩𐨎 𐨬𐨡𐩐 𐨢𐨪𐨿𐨨𐨎 𐨕𐨪𐩐 𐨯𐨿𐨬𐨌𐨢𐨿𐨩𐨌𐨩𐨌𐨣𐨿 𐨨𐨌 𐨤𐨿𐨪𐨨𐨡𐨏𐩐\par} \\
Thai &
  \TH{เมื่อเปิดหนังสือเก่าในห้องสมุด เราอาจพบลายมือของผู้อ่านจากหลายยุคสมัย บางคนขีดเส้นใต้ข้อความที่ชอบ บางคนเขียนคำถามไว้ริมหน้า และบางคนบันทึกวันเวลาที่อ่านจบ เรื่องราวเหล่านี้ทำให้หนังสือเล่มเดียวเป็นเหมือนบทสนทนาที่ดำเนินต่อไปข้ามกาลเวลา} \\
\end{TableBookX}

\noindent
The Segoe UI Historic face is a user-supplied Microsoft font; its
shaping may differ from the redistributable Noto Sans Brahmi.
\index{Brahmi}\index{Kharoshthi}\index{Thai line breaking}

\subsection{Egyptian Hieroglyphic Composition}

The opening of the \emph{Tale of the Shipwrecked Sailor} supplies an
attested sequence of Egyptian quadrats, not an invented arrangement.
The first three numbered lines below follow Bob Richmond's horizontal
UMdC transcription, adapted from JSesh exemplars with reference to
James P. Allen's edition. Editorial colour and line markup have been
omitted; the original EHFC joining and insertion characters remain.
\index{Egyptian Hieroglyph Format Controls}

\noindent\textit{Line 1}\par
\noindent{\large\EG{𓆓𓐳𓂧𓇋𓐰𓈖𓌞𓅱𓀀𓇋𓈎𓐰𓂋𓐰𓏛𓅱𓍒𓄿𓏜}}\par
\medskip
\noindent\textit{Line 2}\par
\noindent{\large\EG{𓄣𓐱𓏤𓐰𓎡𓄂𓐰𓂝𓀀𓅓𓂝𓐰𓎡𓄖𓐰𓂻𓈖𓐰𓈖𓐰𓏥}}\par
\medskip
\noindent\textit{Line 3}\par
\noindent{\large\EG{𓄚𓐰𓈖𓏌𓅱𓉐𓊏𓐰𓊪𓂡𓐍𓐰𓂋𓊪𓐱𓏲𓐰𓆱}}\par

\smallskip
\noindent
Source: \href{https://github.com/HieroglyphsEverywhere/UMdC/blob/master/Samples/ShipwreckedSailorHorizontal.umdc}
{ShipwreckedSailorHorizontal.umdc}. The extant manuscript is in
hieratic script; this is a modern hieroglyphic transcription.

\subsection{Khitan Small Script}

XSR composes Khitan Small Script characters into two-dimensional
clusters directly from their Unicode source. Ordinary clusters and
the Type B form marked by U+16FE4 can be compared without an
intermediate linear rendering.
\index{Khitan Small Script}

\begin{center}
\textbf{Ordinary clusters}\par
\smallskip
\KHS{𘲇 𘰺𘱹𘰻 𘱿𘱤 𘱚𘮒𘲦}

\medskip
\textbf{Type B and additional clusters}\par
\smallskip
\KHS{𘮦𖿤𘲲 𘰕𘲆𘮒 𘲵𘯺𘬾𘯺}
\end{center}


% ==========================================================================
\section{East Asian Writing and CJK Family Selection}
% ==========================================================================

\Term[sort=CJK family selection]{CJK family selection}
[continuous text, shared Han ownership, and explicit regional overrides]
distinguishes the document CJK family from language-specific local families.
Explicit local family selection determines the intended regional form
of shared Han characters, while several East Asian writing systems also
require shaping behavior that extends beyond ordinary Han glyph selection.

The specimens in this section use normal-width continuous text. They exercise paragraph composition,
line breaking, punctuation, regional glyph forms, Hangul composition,
historical jamo, supplementary-plane ideographs, and mixed-script runs.
A compact table at the end isolates shared Han codepoints for direct visual
comparison.


% --------------------------------------------------------------------------
\subsection{Chinese}
% --------------------------------------------------------------------------

\paragraph{Mainland Simplified Chinese --- Noto SC}

This paragraph uses the explicit Simplified Chinese family throughout.

\SC{现代中文排版并不只是把一个个汉字依次放进页面。一个真正的正文环境还需要同时处理标点、数字、书名号、引号、括号以及连续句子之间的换行关系，并且在较长的段落中保持稳定的字面和行距。对于共享的汉字码位，不同地区使用的字体还可能给出不同的字形，因此同一个“骨、直、令、具、角”在不同字体家族中会表现出细微而稳定的区域差异。这里使用一整段连续的简体中文来检查大陆简体字形、标点悬挂附近的布局、正常段落换行以及本地字体选择在长文本中是否能够始终保持有效，而不是只验证几个孤立字符能不能显示。}


\paragraph{Mainland Traditional Chinese --- Noto SC}

Traditional characters can also be deliberately typeset through the Mainland
Chinese family. This separates character repertoire from regional glyph
selection: the explicit family choice supplies the regional form for
traditional codepoints.

\SC{現代中文排版並不只是把一個個漢字依次放進頁面。一個真正的正文環境還需要同時處理標點、數字、書名號、引號、括號以及連續句子之間的換行關係，並且在較長的段落中保持穩定的字面和行距。對於共享的漢字碼位，不同地區使用的字體還可能給出不同的字形，因此同一個「骨、直、令、具、角」在不同字體家族中會表現出細微而穩定的區域差異。這裡特意使用傳統漢字而仍然選擇大陸字體家族，以檢查字符繁簡與地域字形選擇能否保持彼此獨立，並觀察長段落中共享漢字、繁體字和全角標點的連續排版。}


\paragraph{Taiwan Traditional Chinese --- Noto TC}

The Taiwan family is tested with continuous Taiwan Traditional Chinese rather
than with a short regional-glyph sample.

\TC{正體中文的正文排版除了字符本身之外，也牽涉字型、標點、行距與段落節奏。當文件同時包含不同語言或不同地區的漢字時，程式不能僅依 Unicode 碼位猜測應採用哪一套地區字形，而應由明確的語言或字型選擇來決定。這個段落使用臺灣正體字型，並刻意包含「骨、直、令、具、角、門、國、學、龍」等適合觀察地區字形差異的共享漢字，同時加入引號、頓號、逗號與句號，藉此確認較長的正體中文內容在正常欄寬中能自然換行，而且從段首到段末始終維持同一個本地 CJK family。}


\paragraph{Document Traditional Chinese --- Shanggu}

The following paragraph is deliberately left completely unwrapped. It
therefore remains under the document CJK family, Shanggu, and provides the
control against which the explicit Chinese, Japanese, and Korean families can
be compared.

漢字進入現代數位排版之後，同一個編碼往往可以由許多不同的字體呈現，而字形的歷史層次、地域習慣和設計風格並不會因此消失。傳統中文尤其適合用來觀察這種情形：有些字在不同地區的正文中結構略異，有些字的筆形和部件位置不同，而另一些差別只有在一整段文字的灰度與節奏中才會顯現。這個段落不使用任何語言專用命令，使所有漢字直接落入文件的 Shanggu CJK family；因此它不只是檢查單字是否存在，也檢查文件級 CJK 選擇在較長的連續文本、標點和自動換行之間能否穩定維持。


% --------------------------------------------------------------------------
\subsection{Korean and Historical Hangul}
% --------------------------------------------------------------------------

\paragraph{Modern Korean --- Hangul}

The first Korean specimen deliberately avoids Hanja. It tests ordinary modern
Hangul composition, punctuation, word spacing, and paragraph wrapping.

\KR{한글로 쓰인 긴 문장은 여러 음절 블록이 끊임없이 이어지면서도 낱말 사이의 간격과 문장 부호의 위치가 자연스럽게 유지되어야 한다. 짧은 보기에서는 글꼴이 몇 글자를 제대로 보여 주는지만 확인할 수 있지만, 실제 문단에서는 줄 끝에서 낱말이 어떻게 배치되는지, 괄호와 따옴표가 앞뒤 글자와 어떻게 어울리는지, 여러 줄을 지나도 같은 글꼴 선택이 계속 유지되는지를 함께 살펴볼 수 있다. 이 문단은 완성형 한글이 이어지는 일반적인 한국어 본문을 충분히 길게 배치하여 문장 전체의 모양과 줄 간격을 확인하고, 뒤에 나오는 한자 혼용문과 옛한글 자료를 비교하기 위한 현대 한글 기준 표본으로 사용한다.}


\paragraph{Modern Korean --- mixed Hanja and Hangul}

The next paragraph keeps Hangul and Hanja inside a single explicit Korean
family. Shared Han characters retain Korean regional forms throughout the
surrounding Hangul text.

\KR{韓國語의 現代 文體는 오늘날 대체로 한글만으로 적지만, 歷史ㆍ言語ㆍ哲學ㆍ法律과 같은 분야에서는 漢字를 함께 적은 문헌을 어렵지 않게 만날 수 있다. 이런 漢字混用文을 조판할 때에는 한글 音節과 句讀點 사이에 들어가는 漢字가 갑자기 다른 地域 字形으로 바뀌지 않아야 하며, 韓國ㆍ朝鮮ㆍ漢字ㆍ文化ㆍ文字와 같은 어휘도 앞뒤의 한글 문맥과 한 family 안에서 자연스럽게 이어져야 한다. 이 문단은 國語의 읽기 흐름을 유지하면서 여러 漢字語와 고유어를 섞어, 실제 문헌에서 볼 수 있는 긴 혼용문이 줄바꿈을 거쳐도 안정적으로 조판되는지를 시험한다.}


\paragraph{Jeju language --- corpus specimen}

Jeju is useful here because practical orthographies make active use of the
arae-a and decomposed Hangul sequences that are rare in ordinary modern
Korean. The following is a composite regression specimen assembled from
short attested Jeju-language passages for a paragraph-level regression.

\KR{제주말(濟州말)은 제주섬서 ᄀᆞᆯ아지는 말이다. 제주말모음은 시간 지나멍 여러 가지 변화가 이섯다. 제주섬은 한라산광 ᄌᆞᄁᆞᆺ듸 오름덜록 일뤄진 섬이라. 엿적 화산활동으록 굼부리가 꿰쉉 곤데곤데 ᄋᆢ라 오름이 하영 삿다. 낭(木)은 나무질로 뒌 줄기를 가지고 잇인 여러헤살이 식물이다. 제주말 위키벡과 또시는 제주어 위키백과는 제주말로 운영헴신 위키벡과 언어판이라.}

\paragraph{Middle Korean and Old Hangul --- historical corpus specimen}

The historical specimen deliberately retains decomposed jamo and tone marks.
Short genuine historical fragments retain their source spellings and are
placed consecutively to form a paragraph-level
shaping regression.

\KR{나랏〮말〯ᄊᆞ미〮中듀ᇰ國귁〮에〮 달아〮문ᄍᆞᆼ와〮로서르ᄉᆞᄆᆞᆺ디〮아니〮ᄒᆞᆯᄊᆡ〮　아〮기 ᄇᆡ여〮셔 과ᄀᆞᆯ이〮 가ᄉᆞᆷ〮 알파〮주글〮 ᄃᆞᆺ〮ᄒᆞ〮거든〮　두들겐 開闢ᄒᆞᆯ 젯 므리 疏通ᄒᆞ고 남ᄀᆞᆫ 古今엣 남기 섯것도다　머리〮 알ᄑᆞ고〮 고〮히 막고〮 머리〮와눈〮괘〮 환티〮 아니〮커든〮　긔〮운이〮 우흐〮로 올아〮 머리〮 알포〮미 ᄣᆞ리ᄂᆞᆫ ᄃᆞᆺ〮거든〮　넙고 큰 天造 ᄉᆞ이예 다ᄉᆞᆯ며 어즈러우민ᄃᆞᆯ 엇디 덛덛ᄒᆞᆫ 數ㅣ리오}

The corpus paragraph above tests realistic historical material. The following
synthetic boundary sequence has a different purpose: it explicitly crosses
Hangul Jamo Extended-A, ordinary Jamo, and Extended-B while remaining inside
one shaping run.

\begin{center}
\Huge
\KR{ꥠᅡퟋ〮}
\end{center}

Modern decomposed control:

\begin{center}
\Huge
\KR{각}
\end{center}


% --------------------------------------------------------------------------
\subsection{Japanese and Ryukyuan}
% --------------------------------------------------------------------------

\paragraph{Modern Japanese}

Modern Japanese combines kana, regional Han forms, punctuation, and frequent
script changes inside the same sentence. The entire paragraph remains under
the explicit Japanese family.

\JP{現代日本語の文章では、漢字、ひらがな、カタカナが一つの段落の中で何度も切り替わり、句読点や括弧、かぎ括弧も連続して現れる。短い文字列だけではフォントが表示できることしか分からないが、長い本文にすると、漢字と仮名の大きさの釣り合い、約物の位置、行末での折り返し、そして次の行へ移った後の字形の連続性まで確認できる。この段落では「日本語」「漢字」「文字」「学校」「文化」「骨」「直」「令」「具」「角」などを通常の文脈に入れ、日本語用の地域字形を保ったまま複数行の文章が安定して組版されるかを試している。}


\paragraph{Historical kana orthography}

The opening sentence below is historical text; the continuation is a
showcase-written regression passage that deliberately retains historical kana
spellings such as \JP{ゐる}, \JP{やう}, and \JP{言ふ}. This keeps the
specimen long enough for paragraph testing while identifying the continuation
as showcase-written text.

\JP{真夏の宿屋は空虚であつた。窓の外には白い雲がゆつくり流れ、庭の木立では蟬が絶えず鳴いてゐる。旅人は古い椅子に腰をおろし、まだ來ない馬車の音を待ちながら、時計の針ばかりを眺めてゐた。誰も何も言はないので、午後の光だけが廊下を少しづつ移つて行くやうに見えた。かうして長い時間が過ぎても、玄關の戸は一度も開かず、遠くの道から聞えて來る物音だけが、時々靜かな宿の空氣を動かしてゐるのであつた。}

\paragraph{Japanese Kanbun}

The first sentence is taken from Rai San'yō's \emph{Nihon Gaishi}; the
remainder is a showcase-written Kanbun-style regression continuation. The
purpose here is typographic: Japanese Han glyphs must remain selected even
throughout the all-Han run.

\JP{外史氏曰、吾讀舊志、見鳥羽帝時數下制符。蓋日本文章之傳、漢籍之學與國史之記相因而久矣。昔之學者讀經史、作序記、論人物、述治亂、多用漢文而各成其體。今取其字列於一段、所以驗日本字形之連續、句讀之位置、長行之折返、以及共享漢字在無假名介入時亦不離其本地字體也。文字雖同出一碼、而骨直令具角門國學龍諸字之形、隨地域字體而微異，故尤宜以連續文章觀之。}


\paragraph{Okinawan --- corpus specimen}

Okinawan is represented by a composite specimen from the Wikimedia Okinawan
test project. Its mixed use of kana and Han characters makes it a useful
additional Japanese-family regression. The orthography reflects this composite
corpus specimen.

\JP{沖縄口ぇー琉球諸語ぬてぃ壱ちやいびーん。うちなーぐちぬ首里那覇方言ぉー琉球国ぬじょーま言葉やびたん。琉球王國や、昔ジジ中心んかい存在さる王国やたん。王や首里天加那志とぅ言ん。首里城や、沖縄県ぬ城やん。琉球王朝ぬ王城っし、沖縄県ぬいっちん大さるぐしくやん。漢字ぇー、大昔唐ぬ国んかいそーじー現りーる文字ぬてぃーちやん。やまとぅ、高麗、るーちゅー、ベトナムぉー、大昔唐ぬ国から漢字伝わてぃちゃん。}

\paragraph{Ainu}

Ainu written in the modern katakana orthography makes extensive use of
Katakana Phonetic Extensions for syllable-final consonants. The following
specimen therefore exercises ordinary katakana together with characters such
as \JP{ㇰ}, \JP{ㇱ}, \JP{ㇺ}, \JP{ㇻ}, and the combining form \JP{ㇷ゚} inside
continuous text.

\JP{イタㇰプリカ　ソモネコㇿカ　シサムモシㇼ　モシㇼソカワ　チヌㇺケニㇱパ　チヌㇺケ　カッケマッ　ウタペラリワ　オカウㇱケタ　クニネネワ　アイヌイタㇰアニ　クイタㇰルウェ　ネワネヤクン　ラモッシワノ　クヤイライケㇷ゚　ネルウェタパンナ。エエパキタ　カニアナㇰネ　アイヌモシㇼ　シシㇼムカ　ニㇷ゚タニコタン　コアパマカ　クネルウェネ　ウッチケクニㇷ゚　ラカサㇰペ。}

% --------------------------------------------------------------------------
\subsection{Vietnamese}
% --------------------------------------------------------------------------

The same constitutional passage is used throughout this subsection so that
the specimens differ primarily in writing system or font selection rather
than in linguistic content. Quốc Ngữ is first typeset with the document Latin
family, Libertinus, and then with the dedicated Vietnamese family. The
corresponding Hán-Nôm text uses the VinaWiki repertoire.

\paragraph{Quốc Ngữ --- document Latin family}

The first specimen is deliberately unwrapped and therefore remains in the
document Latin family, Libertinus. It tests Vietnamese diacritics through
that document-wide font choice.

Trải qua mấy nghìn năm lịch sử, Nhân dân Việt Nam lao động cần cù, sáng tạo, đấu tranh anh dũng để dựng nước và giữ nước, đã hun đúc nên truyền thống yêu nước, đoàn kết, nhân nghĩa, kiên cường, bất khuất và xây dựng nên nền văn hiến Việt Nam. Từ năm 1930, dưới sự lãnh đạo của Đảng Cộng sản Việt Nam do Chủ tịch Hồ Chí Minh sáng lập và rèn luyện, Nhân dân ta tiến hành cuộc đấu tranh lâu dài, đầy gian khổ, hy sinh vì độc lập, tự do của dân tộc, vì hạnh phúc của Nhân dân. Cách mạng tháng Tám thành công, ngày 2 tháng 9 năm 1945, Chủ tịch Hồ Chí Minh đọc Tuyên ngôn độc lập, khai sinh ra nước Việt Nam dân chủ cộng hòa, nay là Cộng hòa xã hội chủ nghĩa Việt Nam.


\paragraph{Quốc Ngữ --- Vietnamese family}

The same text is repeated through the explicit Vietnamese local family,
allowing the dedicated Vietnamese font to be compared directly with
Libertinus under identical paragraph conditions.

\VI{Trải qua mấy nghìn năm lịch sử, Nhân dân Việt Nam lao động cần cù, sáng tạo, đấu tranh anh dũng để dựng nước và giữ nước, đã hun đúc nên truyền thống yêu nước, đoàn kết, nhân nghĩa, kiên cường, bất khuất và xây dựng nên nền văn hiến Việt Nam. Từ năm 1930, dưới sự lãnh đạo của Đảng Cộng sản Việt Nam do Chủ tịch Hồ Chí Minh sáng lập và rèn luyện, Nhân dân ta tiến hành cuộc đấu tranh lâu dài, đầy gian khổ, hy sinh vì độc lập, tự do của dân tộc, vì hạnh phúc của Nhân dân. Cách mạng tháng Tám thành công, ngày 2 tháng 9 năm 1945, Chủ tịch Hồ Chí Minh đọc Tuyên ngôn độc lập, khai sinh ra nước Việt Nam dân chủ cộng hòa, nay là Cộng hòa xã hội chủ nghĩa Việt Nam.}


\paragraph{Hán-Nôm --- VinaWiki repertoire}

The Hán-Nôm specimen corresponds to the Quốc Ngữ passage above. Because the local family uses the
VinaWiki font, the specimen is taken from the VinaWiki transcription itself
and therefore tests the repertoire for which that font is intended.

\HN{𱱇過𠇍𠦳𢆥歷史、人民越南勞動勤劬、創造、鬥爭英勇抵𥩯渃吧𡨺渃、㐌熏𤒘𢧚傳統𢞅渃、團結、仁義、堅強、不屈吧𡏦𥩯𢧚𡋂文獻越南。自𢆥1930、𨑜事領導𧵑黨共產越南由主席胡志明創立吧㷙鍊、人民些進行局鬥爭𥹰󠄁𨱽、𣹓艱苦、犧牲爲獨立、自由𧵑民族、爲幸福𧵑人民。革命𣎃𠔭成功、𣈜2𣎃9𢆥1945、主席胡志明讀宣言獨立、開生𫥨渃越南民主共和、𫢩羅共和社會主義越南。}


% --------------------------------------------------------------------------
\subsection{Tangut}
% --------------------------------------------------------------------------

Wikisource provides a sizeable Unicode Tangut corpus, including poetry,
Buddhist literature, maxims, and the Tangut text of the Liángzhōu bilingual
stele. The following paragraph is therefore a composite corpus specimen:
several short attested sequences from independent Tangut texts are joined
for visual regression. They form a composite of separate historical works.

\TG{𘘞𗳝𘛮𗫉𗥦𗤄，𗼑𗼑𗫉𗰜𗰓𗨛，𘚜𗨤𗤮𘁾𗣃𗅋𘃪，𗼇𗗙𘋥𘝿𘎪，𗵒𗹶𗲈𗃵𘓺𘋨𗶠，𗦵𘇇𗱢𗆐𗲌𗽓𘏋，𗌜𗎫𗠒𗂥𗅋𗚛𘓁𘟂𗫶，𘚶𗢵𘄩𗘴𘜑𗃱𗀸𗀸𗏹𗅋𗁿，𗒘𘓷𗰜𘕿𗅋𗯗𘓁𘟂𗫶，𗌮𗈦𗪙𘝇𗤁𗵘𘆗𘚢，𘈒𗤄𗢳𘒣𗎫𘟀𗢳𗵆，𘘚𘒣𗰜𗎫𗮔𗩱𗫂𘟀𘟂，𗮔𗰜𗋕𗡶𘉐𗭍𗅋𗋃𗫂𗎫𘟂。}


% --------------------------------------------------------------------------
\subsection{WenJin Multi-Plane Coverage}
% --------------------------------------------------------------------------

WenJin is different from the language-specific families above. It is exposed
as one public local family whose internal fallback configuration spans
multiple CJK planes. The paragraph below is therefore intentionally a
coverage stress test: ordinary Han text is
interleaved with supplementary-plane characters while the public command
remains unchanged.

\WJ{中國文字的編碼範圍遠遠超過基本多文種平面，因此只測試一二三天地玄黃這類常用字並不足以確認大型漢字字體的覆蓋是否正常。這一段先由普通漢字開始，隨後插入第二平面的𠀀、𠮷、𡈽、𡌛、𣎴，再回到典籍文字與古今學術，之後繼續進入第三平面的𰀀、𰀁、𰀂、𰀃、𰀄。當不同平面的字符在同一個長段落中交錯出現時，使用者仍然只看見一個 \texttt{\textbackslash WJ} family；Plane 0、Plane 2 和 Plane 3 的具體 fallback 應由內部配置完成，而不應要求正文作者為不同 Unicode plane 手動切換命令。}

The following table complements the paragraph test with deterministic block
coverage samples. For each CJK Unified Ideographs block, the first twenty
encoded ideographs are shown. This makes gaps in coverage or an incorrect
Plane 0, Plane 2, or Plane 3 fallback immediately visible and keeps the
regression sample stable across releases.

\begin{TableBookX}
  {L L L}
  {WenJin CJK block coverage}
  {tab:wenjin-blocks}

Block & Range & Twenty-character specimen \\
\midrule

CJK Unified Ideographs
& U+4E00--U+9FFF
& \WJ{一丁丂七丄丅丆万丈三上下丌不与丏丐丑丒专} \\

Extension A
& U+3400--U+4DBF
& \WJ{㐀㐁㐂㐃㐄㐅㐆㐇㐈㐉㐊㐋㐌㐍㐎㐏㐐㐑㐒㐓} \\

Extension B
& U+20000--U+2A6DF
& \WJ{𠀀𠀁𠀂𠀃𠀄𠀅𠀆𠀇𠀈𠀉𠀊𠀋𠀌𠀍𠀎𠀏𠀐𠀑𠀒𠀓} \\

Extension C
& U+2A700--U+2B73F
& \WJ{𪜀𪜁𪜂𪜃𪜄𪜅𪜆𪜇𪜈𪜉𪜊𪜋𪜌𪜍𪜎𪜏𪜐𪜑𪜒𪜓} \\

Extension D
& U+2B740--U+2B81F
& \WJ{𫝀𫝁𫝂𫝃𫝄𫝅𫝆𫝇𫝈𫝉𫝊𫝋𫝌𫝍𫝎𫝏𫝐𫝑𫝒𫝓} \\

Extension E
& U+2B820--U+2CEAF
& \WJ{𫠠𫠡𫠢𫠣𫠤𫠥𫠦𫠧𫠨𫠩𫠪𫠫𫠬𫠭𫠮𫠯𫠰𫠱𫠲𫠳} \\

Extension F
& U+2CEB0--U+2EBEF
& \WJ{𬺰𬺱𬺲𬺳𬺴𬺵𬺶𬺷𬺸𬺹𬺺𬺻𬺼𬺽𬺾𬺿𬻀𬻁𬻂𬻃} \\

Extension I
& U+2EBF0--U+2EE5F
& \WJ{𮯰𮯱𮯲𮯳𮯴𮯵𮯶𮯷𮯸𮯹𮯺𮯻𮯼𮯽𮯾𮯿𮰀𮰁𮰂𮰃} \\

Extension G
& U+30000--U+3134F
& \WJ{𰀀𰀁𰀂𰀃𰀄𰀅𰀆𰀇𰀈𰀉𰀊𰀋𰀌𰀍𰀎𰀏𰀐𰀑𰀒𰀓} \\

Extension H
& U+31350--U+323AF
& \WJ{𱍐𱍑𱍒𱍓𱍔𱍕𱍖𱍗𱍘𱍙𱍚𱍛𱍜𱍝𱍞𱍟𱍠𱍡𱍢𱍣} \\

Extension J
& U+323B0--U+3347F
& \WJ{𲎰𲎱𲎲𲎳𲎴𲎵𲎶𲎷𲎸𲎹𲎺𲎻𲎼𲎽𲎾𲎿𲏀𲏁𲏂𲏃} \\

\end{TableBookX}


% --------------------------------------------------------------------------
\subsection{Regional Forms of Shared Han Characters}
% --------------------------------------------------------------------------

The preceding specimens test normal paragraphs. Here shared
codepoints deliberately isolated. Every cell below contains the same Unicode
character; the visible difference should therefore come from the selected
regional family selected for the common codepoint.

\begin{TableBookX}
  {L C C C C C}
  {Regional forms of shared Han codepoints}
  {tab:cjk-regional-forms}

Character & Shanggu & Noto SC & Noto TC & Japanese & Korean \\
\midrule

骨 & 骨 & \SC{骨} & \TC{骨} & \JP{骨} & \KR{骨} \\

直 & 直 & \SC{直} & \TC{直} & \JP{直} & \KR{直} \\

令 & 令 & \SC{令} & \TC{令} & \JP{令} & \KR{令} \\

具 & 具 & \SC{具} & \TC{具} & \JP{具} & \KR{具} \\

角 & 角 & \SC{角} & \TC{角} & \JP{角} & \KR{角} \\

門 & 門 & \SC{門} & \TC{門} & \JP{門} & \KR{門} \\

國 & 國 & \SC{國} & \TC{國} & \JP{國} & \KR{國} \\

學 & 學 & \SC{學} & \TC{學} & \JP{學} & \KR{學} \\

龍 & 龍 & \SC{龍} & \TC{龍} & \JP{龍} & \KR{龍} \\

\end{TableBookX}

% ==========================================================================
\section{Document Features}
% ==========================================================================

\Term[sort=feature system]{feature system}
[composable document modules] is demonstrated throughout this PDF, but this
section gives each major public feature an explicit specimen as well. The
showcase therefore tests individual modules and their coexistence in a
realistic document.

% --------------------------------------------------------------------------
\subsection{Mathematics}
% --------------------------------------------------------------------------

Mathematical font selection is deliberately left at IMPE's default setting.
The following single display combines several common mathematical alphabets,
operators, delimiters, accents, matrices, integrals, sums, limits, radicals,
and nested scripts as a compact stress specimen.

\begingroup
\small
\[
\begin{aligned}
\mathcal Z_{\beta,\mu}
&\coloneqq
\operatorname{Tr}_{\mathscr H}
\exp\!\left[
-\beta
\left(
\widehat{\mathbf H}
-\mu\widehat{\mathbf N}
\right)
\right]
\\[0.5em]
&=
\sum_{n=0}^{\infty}
\frac{1}{n!}
\int_{\mathbb R^{3n}}
\prod_{j=1}^{n}
\frac{d^3\mathbf p_j}{(2\pi\hbar)^3}
\,
\exp\!\left[
-\sum_{j=1}^{n}
\sqrt{
m_j^2c^4+c^2\lVert\mathbf p_j\rVert^2
}
\right]
\\[0.7em]
&\quad\times
\det
\begin{pmatrix}
\displaystyle
\frac{\partial^2\mathcal L}
     {\partial q^i\partial q^j}
&
\displaystyle
\nabla_\mu F^{\mu\nu}
&
\displaystyle
\left\langle
\psi
\middle|
\widehat{\mathbf H}
\middle|
\phi
\right\rangle
\\[0.9em]
\displaystyle
\sum_{k=1}^{\infty}\frac{(-1)^{k+1}}{k^2}
&
\displaystyle
\sqrt[3]{1+\sqrt{1+\lambda^2}}
&
\displaystyle
\frac{\mathbb E[X^4]-3\mathbb E[X^2]^2}
     {\operatorname{Var}(X)^2}
\\[0.9em]
\displaystyle
\oint_{|z|=1}
\frac{\log(1+z)}{z^{r+1}}\,dz
&
\displaystyle
\prod_{p\in\mathbb P}
\left(1-p^{-s}\right)^{-1}
&
\displaystyle
\lim_{\varepsilon\to0^+}
\frac{f(x+\varepsilon)-f(x-\varepsilon)}
     {2\varepsilon}
\end{pmatrix}
\\[1em]
&\quad+
\underbrace{
\widetilde{\mathbf A}^{\,\dagger}{}_{\mu}{}^{\nu}
\,
\overline{\mathbf B}_{\nu}{}^{\rho}
\,
\widehat{\mathbf C}_{\rho}{}^{\mu}
}_{\text{accented tensor contraction}}
+
\frac{
\Gamma\!\left(\frac12+i\alpha\right)
\zeta(2s)
}{
\sqrt{2\pi\sigma^2}
}
\\[0.8em]
&\quad+
\begin{cases}
\displaystyle
e^{-1/x^2}\sin(\pi/x),
& x\ne0,
\\[0.8em]
0,
& x=0,
\end{cases}
\\[0.8em]
&\quad
\xrightarrow[\alpha\to0]{\beta\to\infty}
\left(
\bigoplus_{j=1}^{r}\mathbb C^{\,n_j}
\right)
\otimes
\bigwedge\nolimits^{k}V
\\[0.8em]
&\quad
\Longrightarrow
\mathbb P
\left(
|X-\mathbb E[X]|\ge t
\right)
\le
\frac{\operatorname{Var}(X)}{t^2},
\qquad
t>0,
\\[0.8em]
&\quad
\alpha,\beta,\gamma,\delta,\varepsilon,\theta,\vartheta,\lambda,\mu,\nu,
\xi,\pi,\varpi,\rho,\varrho,\sigma,\phi,\varphi,\chi,\psi,\omega
\in\mathbb R .
\end{aligned}
\]
\endgroup

\subsection{Tables}

Table~\ref{tab:feature-summary} uses the IMPE book-table environment rather
than a raw \texttt{tabular} wrapper.

\begin{TableBookX}
  {L C C}
  {Feature integration summary}
  {tab:feature-summary}
Feature & Public route & Exercised here \\
\midrule
Mathematics & \texttt{math} & yes \\
Tables & \texttt{tables} & yes \\
Images & \texttt{image} & yes \\
Citations & \texttt{citations} & yes \\
Index & \texttt{index} & yes \\
Hyperlinks & \texttt{hyperlinks} & yes \\
Running headers & \texttt{headers} & yes \\
List environments & \texttt{lists\_envs} & yes \\
\end{TableBookX}

\subsection{Lists and Example Blocks}

A normal list remains ordinary \LaTeX:

\begin{enumerate}
  \item automatic font ownership;
  \item script-specific behavior where required;
  \item reusable document features;
  \item stable public interfaces.
\end{enumerate}

The IMPE \texttt{ExampleBlock} environment provides a distinct indented
specimen:

\begin{ExampleBlock}
A canonical showcase should be readable as a document, useful as a visual
regression specimen, and representative of the public interface that new
documents are expected to use.
\end{ExampleBlock}

\subsection{Images}

The image feature is exercised with the standard \TeX~distribution's example
image when it is available. This keeps the canonical source self-contained
using a generated illustration in the build.

\IfFileExists{example-image-a.pdf}{%
  \begin{OneImage}
    [htbp]
    [0.55\linewidth]
    [0.32\textheight]
    {example-image-a}
    [Image-feature specimen using the standard \TeX~example image.]
    [fig:showcase-image]
  \end{OneImage}

  Figure~\ref{fig:showcase-image} also exercises cross-reference hyperlinks.
}{%
  \begin{center}
    \fbox{%
      \parbox{0.62\linewidth}{%
        \centering
        \vspace{2em}
        \textbf{Image specimen unavailable}\\[0.5em]
        \texttt{example-image-a.pdf} was not found in this \TeX~installation.
        \vspace{2em}
      }%
    }
  \end{center}
}

\subsection{Citations and Bibliography}

Citation support is loaded through the feature system and configured before
feature activation. A normal textual citation can therefore be written as
\textcite{knuth1984texbook}, while a parenthetical citation appears as
\parencite{knuth1984texbook}.

The bibliography is printed near the end of the document.

\subsection{Index Terms}

Several concepts in earlier sections were inserted with the public
\texttt{\textbackslash Term} interface. For example,
\Term[sort=canonical showcase]{canonical showcase}
[release-time specimen] itself is now indexed. The generated index at the end
of the PDF therefore also exercises indexed terms, optional descriptions, and
hyperlinks back to their first occurrence.

\subsection{Hyperlinks, Footnotes, and Running Headers}

The table of contents is clickable, numbered headings carry stable targets,
and section headings can link back toward their table-of-contents entries.
The running header above this page is produced by the \texttt{headers}
feature.

Footnote markers are also bidirectional.\footnote{Clicking the footnote marker
jumps to this note; the note marker can link back to the original position.}

For an external-link specimen, the project repository can be represented with
a normal hyperlink:
\href{https://github.com/KelvinYangBK67/IMPE-LaTeX-System}
{IMPE \LaTeX~System on GitHub}.


\printbibliography[
  heading=bibintoc,
  title={References}
]

\printindex

\end{document}
